\documentclass[11pt,letterpaper]{article}

\usepackage[top=1in, bottom=1in, left=1in, right=1in, headheight=14pt]{geometry}
\usepackage{fontspec}
\setmainfont{DejaVu Serif}
\setsansfont{DejaVu Sans}
\setmonofont{DejaVu Sans Mono}

\usepackage{xcolor}
\usepackage{titlesec}
\usepackage{fancyhdr}
\usepackage{enumitem}
\usepackage{hyperref}
\usepackage{booktabs}
\usepackage{tabularx}
\usepackage{longtable}
\usepackage{colortbl}
\usepackage{multirow}
\usepackage{array}
\usepackage{parskip}
\usepackage{tcolorbox}
\usepackage{graphicx}
\usepackage{lastpage}

% === COLOR SCHEME: Technology (Steel Blue / Electric / Graphite) ===
\definecolor{primary}{HTML}{263238}
\definecolor{secondary}{HTML}{1565C0}
\definecolor{tertiary}{HTML}{37474F}
\definecolor{accent}{HTML}{0277BD}
\definecolor{lightprimary}{HTML}{ECEFF1}
\definecolor{lightsecondary}{HTML}{E3F2FD}
\definecolor{lighttertiary}{HTML}{ECEFF1}
\definecolor{rowgray}{HTML}{F5F5F5}
\definecolor{bordergray}{HTML}{BDBDBD}
\definecolor{subtlegray}{HTML}{757575}
\definecolor{darktext}{HTML}{212121}

% === SECTION FORMATTING ===
\titleformat{\section}
  {\Large\bfseries\sffamily\color{primary}}
  {\thesection}{1em}{}
  [\vspace{-0.5em}\textcolor{bordergray}{\rule{\textwidth}{0.4pt}}]

\titleformat{\subsection}
  {\large\bfseries\sffamily\color{secondary}}
  {\thesubsection}{1em}{}

\titleformat{\subsubsection}
  {\normalsize\bfseries\sffamily\color{tertiary}}
  {\thesubsubsection}{1em}{}

\titlespacing*{\section}{0pt}{1.5em}{0.8em}
\titlespacing*{\subsection}{0pt}{1.2em}{0.5em}
\titlespacing*{\subsubsection}{0pt}{0.8em}{0.3em}

% === CUSTOM ENVIRONMENTS ===
\newcommand{\stageheader}[2]{%
  \noindent\colorbox{#2}{\makebox[\textwidth][l]{%
    \textcolor{white}{\bfseries\sffamily\hspace{6pt}#1\hspace{6pt}}}}%
  \vspace{0.5em}\par%
}

\newtcolorbox{asciibox}{
  colback=rowgray, colframe=bordergray,
  arc=1pt, boxrule=0.5pt,
  left=6pt, right=6pt, top=4pt, bottom=4pt,
  fontupper=\small\ttfamily,
  breakable
}

\newtcolorbox{quotebox}{
  colback=lightprimary, colframe=primary,
  arc=2pt, boxrule=0.5pt,
  left=12pt, right=12pt, top=8pt, bottom=8pt,
  breakable
}

\newcommand{\metafield}[2]{\textbf{\sffamily #1} #2\par}
\newcolumntype{L}[1]{>{\raggedright\arraybackslash}p{#1}}
\newcolumntype{C}[1]{>{\centering\arraybackslash}p{#1}}
\newcommand{\tableheader}[1]{\textbf{\sffamily\textcolor{white}{#1}}}

% === HEADERS AND FOOTERS ===
\pagestyle{fancy}
\fancyhf{}
\fancyhead[L]{\small\sffamily\textcolor{subtlegray}{PR \#28 Tessl Review}}
\fancyhead[R]{\small\sffamily\textcolor{subtlegray}{GSD Mission Package}}
\fancyfoot[C]{\small\sffamily\textcolor{subtlegray}{Page \thepage\ of \pageref{LastPage}}}
\renewcommand{\headrulewidth}{0.4pt}
\renewcommand{\footrulewidth}{0pt}

% === HYPERLINKS ===
\hypersetup{
  colorlinks=true,
  linkcolor=secondary,
  urlcolor=secondary,
  pdfauthor={GSD Ecosystem},
  pdftitle={PR 28 Tessl Skill Review -- Mission Package},
  pdfsubject={Vision to Mission Pipeline},
}

\begin{document}

% ================================================================
% TITLE PAGE
% ================================================================
\thispagestyle{empty}
\begin{center}
\vspace*{3cm}

{\small\sffamily\textcolor{primary}{\textbf{GSD RESEARCH MISSION PACKAGE}}}

\vspace{0.3cm}
\textcolor{bordergray}{\rule{8cm}{0.4pt}}

\vspace{1.5cm}
{\fontsize{28}{34}\selectfont\bfseries\sffamily\textcolor{primary}{Tessl Skill Review\\[0.3em]PR \#28 Analysis}}

\vspace{0.8cm}
{\Large\sffamily\textcolor{secondary}{Agent Skills Ecosystem \& GSD Architectural Alignment}}

\vspace{0.5cm}
{\large\sffamily\itshape\textcolor{tertiary}{A Deep Research Study}}

\vspace{3cm}
{\sffamily\textcolor{accent}{Vision $\rightarrow$ Research $\rightarrow$ Mission}}\\[0.3em]
{\small\sffamily\textcolor{accent}{Full Pipeline Execution}}

\vspace{3cm}
{\small\sffamily\textcolor{subtlegray}{Date: March 26, 2026  |  Status: Ready for GSD Orchestrator}}\\[0.3em]
{\small\sffamily\textcolor{subtlegray}{Activation Profile: Squadron  |  Estimated: 8 context windows across 3 sessions}}
\end{center}
\newpage

% ================================================================
% TABLE OF CONTENTS
% ================================================================
\tableofcontents
\newpage

% ================================================================
% STAGE 1 — VISION DOCUMENT
% ================================================================
\stageheader{STAGE 1 --- VISION DOCUMENT}{primary}

\section{PR \#28 Tessl Skill Review --- Vision Guide}

\metafield{Date:}{March 26, 2026}
\metafield{Status:}{Active Research / PR Review Decision Pending}
\metafield{Depends on:}{gsd-skill-creator-analysis.md, gsd-chipset-architecture-vision.md, gsd-mathematical-foundations-conversation.md}
\metafield{Context:}{Evaluating an unsolicited pull request from a Tessl employee that proposes modifications to five GSD skill-creator skills plus installation of a third-party GitHub Action workflow. Requires understanding Tessl's platform, the Agent Skills specification, the specific changes proposed, security implications, and alignment with GSD's architectural philosophy before making a merge/reject/adapt decision.}

\subsection{Vision}

On March 26, 2026, a Tessl employee (rohan-tessl) submitted PR \#28 to the gsd-skill-creator repository, targeting the dev branch. The PR modifies five SKILL.md files in the mathematical foundations engine (MFE) domain plus the DACP interpreter skill, and adds a GitHub Action workflow that integrates the Tessl CLI into the project's CI pipeline.

This is not a hostile action --- Tessl is a legitimate, well-funded (\$125M raised) company founded by Guy Podjarny (Snyk founder, ex-Akamai CTO) that builds a package manager and evaluation platform for agent skills. The PR follows conventional open-source contribution patterns: fork, branch from dev, conventional commits, honest disclosure of employer affiliation. The changes are narrowly scoped and the PR description is thorough.

However, the PR touches skills that encode the mathematical and philosophical foundations of ``The Space Between'' --- a 923-page mathematical textbook that serves as the intellectual spine of the entire GSD ecosystem. These are not generic utility skills. The MFE domain skills encode specific pedagogical and epistemological structures: the Complex Plane of Experience, domain positioning via plane coordinates and radii, and a carefully designed progression from perception to generation. Changes that score well on a generic ``skill quality'' rubric may actively damage the semantic architecture these skills carry.

The Amiga Principle applies here directly: the power of these skills is not in their superficial structure, but in the \emph{architectural relationships between them} --- the spaces between. A tool that scores structural compliance against a generic specification cannot evaluate whether the inter-skill relationships, the mathematical positioning metadata, or the pedagogical through-line have been preserved.

\subsection{Problem Statement}

\begin{enumerate}[leftmargin=2em]
\item \textbf{Context-blind optimization.} The Tessl review tool scores skills against the agentskills.io specification --- a generic structural standard. The GSD MFE skills encode domain-specific mathematical architecture (plane coordinates, domain radii, chapter mappings, primitive counts) that the spec knows nothing about. High Tessl scores and high GSD utility are orthogonal metrics.

\item \textbf{Summary section flattening.} The PR replaces structured metadata in the Summary sections (separate fields for Part name, Chapters, Plane Position, Primitives, Key Concepts) with single-line compressed formats. This makes the information less parseable by agents that consume these fields programmatically and collapses the visual hierarchy that supports human reading during development.

\item \textbf{Workflow section injection.} The PR adds 5-step ``Workflow'' sections to MFE domain skills that were not designed to have procedural workflows. These are \emph{reference} skills (mathematical domains with concept definitions), not \emph{procedural} skills (step-by-step task execution). Adding procedural workflows to reference skills misrepresents their activation pattern.

\item \textbf{Third-party CI integration.} The GitHub Action installs the Tessl CLI and runs it on every PR that touches SKILL.md files. While it uses only GITHUB\_TOKEN (no Tessl account required), it introduces a dependency on \texttt{tesslio/skill-review@main} --- an unpinned reference to a third-party action that could change behavior at any time. This is a supply-chain consideration.

\item \textbf{DACP interpreter restructuring.} The DACP skill changes are more defensible --- restructuring numbered steps into named sections with an Error Handling section --- but the PR also adds YAML frontmatter with a ``name'' and ``description'' field that follows the agentskills.io spec pattern rather than the GSD skill-creator's existing frontmatter schema (which already has these fields in a different format).
\end{enumerate}

\subsection{Core Concept}

\begin{quotebox}
\textbf{Model:} Evaluate PR \#28 by separating \emph{structurally valid improvements} from \emph{context-blind damage}, extract what's useful, reject what's harmful, and establish governance patterns for future external skill contributions.
\end{quotebox}

\subsubsection{Domain Map}

\begin{asciibox}
\begin{verbatim}
         TESSL ECOSYSTEM                    GSD ECOSYSTEM
    +---------------------+          +------------------------+
    | Tessl CLI            |          | skill-creator          |
    | - skill review       |          | - 6 core functions     |
    | - skill lint         |          | - MFE domain skills    |
    | - skill publish      |          | - DACP interpreter     |
    | - skill eval         |          | - Chipset architecture |
    +----------+----------+          +----------+-------------+
               |                                |
    +----------v----------+          +----------v-------------+
    | Agent Skills Spec    |<-------->| Agent Skills Spec      |
    | (agentskills.io)     |  shared  | (extended w/ GSD       |
    | - YAML frontmatter   |  spec    |   fields: triggers,    |
    | - SKILL.md body      |          |   extends, version,    |
    | - progressive disc.  |          |   enabled, metadata)   |
    +---------------------+          +------------------------+
\end{verbatim}
\end{asciibox}

\subsubsection{Module Architecture}

\begin{asciibox}
\begin{verbatim}
Module 1: TESSL PLATFORM     Module 2: AGENT SKILLS SPEC
  - Company & funding           - Spec structure
  - Review methodology          - Scoring criteria
  - Scoring rubric              - GSD extensions vs. spec
  - Similar PR patterns         - Compliance gaps

Module 3: CHANGE IMPACT       Module 4: GOVERNANCE
  - Per-skill diff analysis     - GitHub Action security
  - MFE metadata preservation   - Unpinned action risks
  - DACP restructuring eval     - Contribution policy
  - Workflow appropriateness    - Future PR patterns

Module 5: RESPONSE STRATEGY
  - Accept/reject/adapt matrix
  - Cherry-pick recommendations
  - GSD-native improvements
  - Tessl relationship framing
\end{verbatim}
\end{asciibox}

\subsection{Research Modules}

\subsubsection{Module 1: Tessl Platform Analysis}
Comprehensive analysis of Tessl as a company, its skill review methodology, scoring rubric (validation checks, implementation quality, activation score via LLM-as-judge), and patterns of engagement with open-source skill repositories. Includes examination of the ``good-oss-citizen'' repository and Product Hunt evidence of proactive PR submission as a growth/community strategy.

\subsubsection{Module 2: Agent Skills Specification Audit}
Deep comparison of the agentskills.io specification against GSD's extended skill format. Maps which spec requirements GSD already meets, which GSD intentionally extends beyond, and which genuine compliance gaps exist. Evaluates whether Tessl's scoring rubric captures anything GSD's own quality processes miss.

\subsubsection{Module 3: Per-Skill Change Impact Analysis}
Line-by-line analysis of all six file changes in PR \#28. For each skill: what was changed, what was gained, what was lost, and whether the net effect is positive, negative, or mixed. Special attention to MFE metadata structure (plane coordinates, primitive counts, chapter mappings) and whether the ``Workflow'' sections are semantically appropriate for reference skills.

\subsubsection{Module 4: Governance and Security Review}
Security analysis of the proposed GitHub Action: supply-chain risks of unpinned \texttt{@main} references, permissions requested (\texttt{pull-requests: write, contents: read}), data flow (does \texttt{tessl skill review} send skill content to Tessl servers?), and whether the action's value justifies its integration cost. Includes contribution policy recommendations for handling future external PRs.

\subsubsection{Module 5: Response Strategy}
Decision matrix for each component of the PR: full accept, partial accept with modifications, reject with explanation, or adapt the ideas independently. Produces a concrete action plan including PR review comment text, any GSD-native improvements inspired by the PR, and relationship framing with Tessl as an ecosystem participant.

\subsection{Chipset Configuration}

\begin{asciibox}
\begin{verbatim}
chipset:
  name: "pr28-tessl-review"
  version: "1.0"
  activation: "squadron"
  skills:
    - skill-creator-analysis    # Existing ecosystem knowledge
    - agent-skills-spec-audit   # New: spec compliance analysis
    - pr-diff-analyzer          # New: per-file change evaluation
    - governance-reviewer       # New: CI/CD security analysis
    - response-composer         # New: PR response generation
  agents:
    FLIGHT: { model: opus, role: "merge decision authority" }
    CRAFT-SPEC: { model: opus, role: "spec compliance specialist" }
    CRAFT-MATH: { model: opus, role: "MFE domain guardian" }
    EXEC-1: { model: sonnet, role: "diff analysis executor" }
    EXEC-2: { model: sonnet, role: "governance analysis executor" }
    VERIFY: { model: sonnet, role: "cross-module consistency" }
\end{verbatim}
\end{asciibox}

\subsection{Scope Boundaries}

\textbf{In Scope (v1.0):}
\begin{itemize}[leftmargin=2em]
\item Complete analysis of all six files changed in PR \#28
\item Tessl platform evaluation (methodology, scoring, business model)
\item Agent Skills spec compliance analysis for affected skills
\item GitHub Action security and supply-chain risk assessment
\item Concrete accept/reject/adapt decision for each PR component
\item PR review comment draft for rohan-tessl
\item GSD-native improvement recommendations inspired by the PR
\end{itemize}

\textbf{Out of Scope:}
\begin{itemize}[leftmargin=2em]
\item Full audit of all GSD skills against agentskills.io spec (future milestone)
\item Tessl CLI installation or integration into GSD build pipeline
\item Rewriting the five affected skills from scratch
\item Evaluating Tessl's paid platform offerings
\item Engaging with Tessl's skill registry or publishing GSD skills to it
\end{itemize}

\subsection{Success Criteria}

\begin{enumerate}[leftmargin=2em]
\item A clear accept/reject/adapt decision documented for each of the six files in the PR
\item MFE domain metadata (plane coordinates, radii, chapter numbers, primitive counts) preservation verified in any accepted changes
\item DACP interpreter restructuring evaluated against actual bundle processing workflow
\item GitHub Action security risks quantified with specific mitigation recommendations
\item PR review comment drafted that is technically precise, respectful, and actionable
\item At least three GSD-native improvements identified that the PR inspired but should be implemented differently
\item Agent Skills spec compliance gaps cataloged for the five affected skills (current state vs. spec requirements)
\item Contribution policy recommendations documented for handling future external skill PRs
\item The ``Workflow'' section pattern evaluated: when is it appropriate for GSD skills and when is it not?
\item Tessl's scoring methodology understood well enough to predict scores for any GSD skill
\end{enumerate}

\subsection{Relationship to Other Documents}

\begin{center}
\rowcolors{2}{rowgray}{white}
\begin{tabularx}{\textwidth}{L{4cm}X}
\toprule
\rowcolor{primary}
\tableheader{Document} & \tableheader{Relationship} \\
\midrule
gsd-skill-creator-analysis.md & Architecture context for understanding how skills are structured, loaded, and scored \\
gsd-mathematical-foundations-conversation.md & Philosophical foundation for why MFE domain skill metadata matters \\
gsd-upstream-intelligence-pack-v1.43.md & Upstream monitoring pattern --- Tessl changes could trigger intelligence events \\
gsd-chipset-architecture-vision.md & Chipset YAML references that depend on skill metadata structure \\
unit-circle-skill-creator-synthesis.md & The mathematical spine that MFE skills encode \\
\bottomrule
\end{tabularx}
\end{center}

\subsection{The Through-Line}

The spaces between are where meaning lives. Tessl's scoring rubric measures the \emph{nodes} --- structural compliance, description completeness, workflow presence. But the GSD MFE skills encode the \emph{arrows between nodes} --- the Complex Plane positioning, the domain radii, the chapter-to-concept mappings that make the mathematical foundations engine a coherent navigation system rather than a collection of disconnected reference files.

This PR is a gift wrapped in a misunderstanding. The gift is external validation that some GSD skills have structural gaps worth addressing. The misunderstanding is that a generic scoring rubric can improve domain-specific architectural skills without understanding what those skills encode. The mission is to unwrap the gift carefully --- taking what's genuinely useful while protecting the architecture that makes GSD's mathematical foundations more than the sum of their parts.

\newpage

% ================================================================
% STAGE 2 — RESEARCH REFERENCE
% ================================================================
\stageheader{STAGE 2 --- RESEARCH REFERENCE}{secondary}

\section{PR \#28 Tessl Skill Review --- Research Reference}

\subsection{How to Use This Document}

This reference compiles findings from deep research into the Tessl platform, the Agent Skills specification, and the specific changes proposed in PR \#28. Key source categories:

\begin{itemize}[leftmargin=2em]
\item \textbf{Tessl documentation:} docs.tessl.io, tessl.io/blog, tessl.io/registry
\item \textbf{Agent Skills specification:} agentskills.io/specification
\item \textbf{Anthropic documentation:} platform.claude.com/docs, github.com/anthropics/skills
\item \textbf{GitHub PR data:} github.com/Tibsfox/gsd-skill-creator/pull/28 (diff, commits, files)
\item \textbf{Tessl GitHub organization:} github.com/tesslio (19 public repositories)
\end{itemize}

\subsection{Module 1: Tessl Platform Reference}

\subsubsection{Company Profile}

Tessl was founded by Guy Podjarny (``Guypo''), who previously founded Snyk (scaled to \$8B valuation, 1000+ employees) and Blaze (acquired by Akamai, where Podjarny served as CTO). Tessl raised \$125M total: a \$25M seed led by boldstart and GV (Google Ventures) in April 2024, and a \$100M Series A led by Index Ventures with Accel participation in November 2024. Headquartered in London. The company's thesis is ``AI native software development'' --- moving from code-centric to spec-centric development where humans express intent and AI handles implementation.

\subsubsection{Tessl's Skill Review Methodology}

Tessl evaluates skills across three dimensions:

\begin{center}
\rowcolors{2}{rowgray}{white}
\begin{tabularx}{\textwidth}{L{3cm}XL{2.5cm}}
\toprule
\rowcolor{primary}
\tableheader{Dimension} & \tableheader{What It Measures} & \tableheader{Method} \\
\midrule
Validation & Frontmatter schema, required fields, section layout, line count, metadata & Automated lint \\
Implementation & Instruction clarity, edge case handling, example quality, conciseness & LLM-as-judge \\
Activation & Description quality, trigger terms, ``Use when\ldots'' clause, distinctiveness & LLM-as-judge \\
\bottomrule
\end{tabularx}
\end{center}

The activation score is particularly relevant to PR \#28. Tessl's rubric evaluates descriptions on: trigger\_term\_quality (do natural user phrases appear?), completeness (does it say both \emph{what} and \emph{when}?), and distinctiveness\_conflict\_risk (could it be confused with other skills?). The ``Use when\ldots'' clause is a key signal --- Tessl scores descriptions significantly higher when they include explicit activation triggers.

\subsubsection{Tessl's Open-Source Engagement Pattern}

Evidence from multiple sources indicates Tessl proactively submits PRs to open-source repositories as a community engagement and product awareness strategy. On Product Hunt, a Tessl representative told a user ``I ran all your skills through the Tessl review machinery and sent you a pull request.'' The tesslio GitHub organization includes a repository called ``good-oss-citizen'' and another called ``skill-review-and-optimize'' --- both indicating formalized outreach workflows. This is legitimate open-source participation with a commercial awareness component, not spam. The honest disclosure in PR \#28 (``I work at @tesslio'') follows best practices.

\subsubsection{The tesslio/skill-review GitHub Action}

The action repository has 10 commits, 4 contributors (including Claude and Copilot as AI-assisted contributors), 2 stars, and is written in TypeScript. It installs the Tessl CLI, runs \texttt{tessl skill review} on changed SKILL.md files, and posts a single PR comment (updated on subsequent pushes, not duplicated). The action requires \texttt{pull-requests: write} and \texttt{contents: read} permissions. It uses only \texttt{GITHUB\_TOKEN} --- no Tessl API token or account is needed.

\textbf{Critical question:} Does \texttt{tessl skill review} run entirely locally or does it phone home? The Tessl documentation states ``runs \texttt{tessl skill review} locally --- no Tessl account or API token needed.'' However, the review includes an ``LLM-as-judge'' component for implementation and activation scoring, which likely requires an API call to an LLM provider. This needs verification before any CI integration.

\subsection{Module 2: Agent Skills Specification Reference}

\subsubsection{Specification Overview}

The Agent Skills specification (agentskills.io) is an open format maintained by Anthropic, adopted by 26+ platforms. A skill is a directory containing a SKILL.md file with YAML frontmatter and Markdown body, plus optional \texttt{scripts/}, \texttt{references/}, and \texttt{assets/} directories. The spec requires only two frontmatter fields: \texttt{name} and \texttt{description}. Optional fields include \texttt{license}, \texttt{compatibility}, \texttt{metadata} (string-to-string map), and \texttt{allowed-tools}.

The spec recommends keeping SKILL.md under 500 lines and instructions under 5,000 tokens. It uses a three-tier progressive disclosure model: metadata loads at startup (30--50 tokens per skill), full SKILL.md loads on activation, and reference files load on demand during execution.

\subsubsection{GSD Extensions Beyond the Spec}

GSD's skill format extends the base spec with several fields not in the agentskills.io standard:

\begin{center}
\rowcolors{2}{rowgray}{white}
\begin{tabularx}{\textwidth}{L{3.5cm}XL{2.5cm}}
\toprule
\rowcolor{primary}
\tableheader{GSD Field} & \tableheader{Purpose} & \tableheader{Spec Status} \\
\midrule
\texttt{user-invocable} & Whether users can directly trigger the skill & GSD extension \\
\texttt{allowed-tools} & Tool access scope & In spec (experimental) \\
\texttt{metadata.*} & Domain-specific key-value pairs (chapter, plane-position, radius, etc.) & Spec allows arbitrary metadata \\
\texttt{triggers} & Activation phrase patterns & GSD extension \\
\texttt{extends} & Skill inheritance chain & GSD extension \\
\texttt{version} & Semantic version & GSD extension \\
\texttt{enabled} & Active/inactive toggle & GSD extension \\
\bottomrule
\end{tabularx}
\end{center}

The MFE domain skills use the \texttt{metadata} field extensively to encode mathematical positioning data (chapter numbers, plane coordinates, domain radii, primitive counts). This is spec-compliant --- the spec explicitly allows arbitrary metadata key-value pairs. Tessl's scoring rubric does not penalize custom metadata, but it also does not evaluate whether custom metadata is preserved during modifications.

\subsubsection{Genuine Compliance Gaps in GSD Skills}

The PR does highlight some real gaps where GSD skills diverge from spec best practices:

\begin{enumerate}[leftmargin=2em]
\item \textbf{Description ``Use when\ldots'' clause:} The four MFE skills had descriptions that stated \emph{what} the domain covers but not \emph{when} to activate it. Since these are non-user-invocable reference skills loaded by the MFE orchestrator (not by direct user request), the ``Use when'' clause matters less --- but it would still help programmatic activation if well-crafted.

\item \textbf{DACP interpreter lacked YAML frontmatter:} The original DACP skill had no \texttt{name} or \texttt{description} in frontmatter format, only a Markdown header and quote block. This is a genuine spec compliance gap.

\item \textbf{No explicit ``Workflow'' or ``Steps'' sections:} The spec does not require these, but Tessl's implementation quality scoring rewards clear procedural structure. For the DACP interpreter (which \emph{is} procedural), this is a valid improvement. For MFE reference skills (which are \emph{not} procedural), injecting workflows misrepresents the skill type.
\end{enumerate}

\subsection{Module 3: Per-Skill Change Impact Reference}

\subsubsection{dacp-interpreter/SKILL.md}

\begin{center}
\rowcolors{2}{rowgray}{white}
\begin{tabularx}{\textwidth}{L{3cm}XX}
\toprule
\rowcolor{primary}
\tableheader{Aspect} & \tableheader{Change} & \tableheader{Assessment} \\
\midrule
Frontmatter & Added YAML \texttt{name} and \texttt{description} & \textbf{Positive} --- genuine spec compliance \\
Summary & Compressed to single paragraph & \textbf{Mixed} --- more concise but lost Safety Invariant callout \\
Steps & Renamed ``Step-by-Step Interpretation'' to numbered ``Step N: Title'' format & \textbf{Positive} --- clearer section navigation \\
Error Handling & Grouped under new ``Error Handling'' section & \textbf{Positive} --- better information architecture \\
Duplicate header & Original had duplicate ``\# DACP Interpreter'' --- PR introduces this bug & \textbf{Negative} --- new bug introduced \\
\bottomrule
\end{tabularx}
\end{center}

\textbf{Net assessment:} Mostly positive. The DACP interpreter \emph{is} a procedural skill, so workflow restructuring is appropriate. However, the PR introduces a duplicate H1 header (the original title plus a new one after the frontmatter) and compresses the Safety Invariant (SAFE-01) callout into the description field, reducing its visibility. The Safety Invariant should remain a prominent, standalone section.

\subsubsection{MFE Domain Skills (change, emergence, synthesis, unification)}

All four MFE skills received identical structural changes:

\begin{center}
\rowcolors{2}{rowgray}{white}
\begin{tabularx}{\textwidth}{L{3cm}XX}
\toprule
\rowcolor{primary}
\tableheader{Aspect} & \tableheader{Change} & \tableheader{Assessment} \\
\midrule
Description & Expanded with ``Use when\ldots'' clause and specific topics & \textbf{Positive} --- better activation signal \\
Summary format & Multi-field format $\rightarrow$ single-line compressed & \textbf{Negative} --- lost structured metadata \\
Key Concepts & Renamed to ``Key Concepts'' (was ``Key Primitives'') & \textbf{Negative} --- ``Primitives'' is the correct GSD term \\
Workflow added & 5-step procedural workflow injected & \textbf{Negative} --- these are reference skills, not procedural \\
Metadata preserved & YAML frontmatter metadata block unchanged & \textbf{Positive} --- mathematical positioning intact \\
\bottomrule
\end{tabularx}
\end{center}

\textbf{Net assessment:} Mixed-to-negative. The description improvements are genuinely useful. But flattening the Summary format, renaming ``Primitives'' to ``Concepts'', and injecting procedural workflows into reference skills damages the semantic architecture. The MFE skills are \emph{knowledge cards} that the MFE orchestrator loads for domain context --- they are not step-by-step procedures an agent follows sequentially.

\textbf{Critical detail:} The YAML frontmatter \texttt{metadata} block (containing chapter numbers, plane position coordinates, domain radius, and primitive counts) was \emph{not} modified by the PR. This is the most architecturally sensitive data and it was preserved. The damage is in the Markdown body, not the machine-readable metadata.

\subsection{Module 4: Governance Reference}

\subsubsection{GitHub Action Supply Chain Risks}

\begin{center}
\rowcolors{2}{rowgray}{white}
\begin{tabularx}{\textwidth}{L{3cm}XL{2cm}}
\toprule
\rowcolor{primary}
\tableheader{Risk} & \tableheader{Description} & \tableheader{Severity} \\
\midrule
Unpinned ref & \texttt{tesslio/skill-review@main} can change without notice & Medium \\
CLI install & Action installs \texttt{tessl} CLI binary at runtime & Medium \\
Network calls & LLM-as-judge likely requires external API call & Unknown \\
Permissions & \texttt{pull-requests: write} allows comment posting & Low \\
Data exposure & Skill content sent to Tessl/LLM for scoring & Unknown \\
\bottomrule
\end{tabularx}
\end{center}

\textbf{Mitigation if adopting:} Pin to a specific commit SHA instead of \texttt{@main}. Verify whether the CLI phones home. Scope to a separate workflow that requires manual approval. Never grant \texttt{contents: write} permission.

\subsubsection{Contribution Policy Gap}

GSD's CONTRIBUTING.md specifies branch-from-dev and conventional commits (both followed by this PR), but does not address: external tool integration in CI, skill content modification by non-domain-experts, or architectural review requirements for skills that encode domain-specific knowledge. This PR reveals the need for a \textbf{skill modification governance policy}.

\subsection{Module 5: Response Strategy Reference}

\subsubsection{Decision Framework}

For each component, evaluate on three axes:

\begin{enumerate}[leftmargin=2em]
\item \textbf{Structural improvement:} Does the change make the skill more compliant with the Agent Skills spec in a way that benefits GSD?
\item \textbf{Semantic preservation:} Does the change preserve the domain-specific meaning, relationships, and architectural intent?
\item \textbf{Integration cost:} Does accepting the change create ongoing maintenance burden or external dependencies?
\end{enumerate}

\subsubsection{Recommended Actions Per Component}

\begin{center}
\rowcolors{2}{rowgray}{white}
\begin{tabularx}{\textwidth}{L{3.5cm}L{2cm}X}
\toprule
\rowcolor{primary}
\tableheader{Component} & \tableheader{Decision} & \tableheader{Rationale} \\
\midrule
DACP frontmatter & \textbf{Accept} & Genuine spec compliance improvement \\
DACP step restructuring & \textbf{Accept} & Better navigation for a procedural skill \\
DACP error handling section & \textbf{Accept} & Improved information architecture \\
MFE description expansion & \textbf{Adapt} & Take the ``Use when\ldots'' pattern, rewrite for GSD context \\
MFE summary flattening & \textbf{Reject} & Destroys structured metadata in Markdown body \\
MFE workflow injection & \textbf{Reject} & Misrepresents reference skills as procedural \\
MFE ``Primitives'' rename & \textbf{Reject} & ``Primitives'' is the correct GSD terminology \\
GitHub Action workflow & \textbf{Reject} & Supply-chain risk, unverified data flow, limited value for GSD's architecture \\
\bottomrule
\end{tabularx}
\end{center}

\subsection{Safety and Sensitivity Considerations}

\begin{center}
\rowcolors{2}{rowgray}{white}
\begin{tabularx}{\textwidth}{L{3cm}XL{2.5cm}}
\toprule
\rowcolor{primary}
\tableheader{Concern} & \tableheader{Details} & \tableheader{Mitigation} \\
\midrule
Professional courtesy & PR author disclosed employer affiliation; response should be respectful and technically grounded & Draft review comment with appreciation + specific technical feedback \\
Supply chain & Third-party action with runtime CLI installation & Do not merge workflow file; evaluate independently if desired \\
IP exposure & Skill content may be sent to external services during review & Verify data flow before any Tessl CLI usage on proprietary skills \\
Mathematical integrity & MFE skills encode ``The Space Between'' foundations & Any modifications must preserve plane coordinates, radii, chapters, and primitive counts \\
\bottomrule
\end{tabularx}
\end{center}

\subsection{Source Bibliography}

\subsubsection{Primary Sources}

\begin{itemize}[leftmargin=2em]
\item GitHub PR \#28: \url{https://github.com/Tibsfox/gsd-skill-creator/pull/28}
\item Agent Skills Specification: \url{https://agentskills.io/specification}
\item Anthropic Agent Skills Documentation: \url{https://platform.claude.com/docs/en/agents-and-tools/agent-skills/overview}
\item Tessl Skill Review Documentation: \url{https://docs.tessl.io/evaluate/evaluating-skills}
\item tesslio/skill-review Action: \url{https://github.com/tesslio/skill-review}
\end{itemize}

\subsubsection{Company and Industry Sources}

\begin{itemize}[leftmargin=2em]
\item Tessl Platform: \url{https://tessl.io/}
\item Tessl Blog --- Skills Announcement: \url{https://tessl.io/blog/skills-are-software-and-they-need-a-lifecycle-introducing-skills-on-tessl/}
\item Tessl Blog --- Three Eval Methodologies: \url{https://tessl.io/blog/three-context-eval-methodologies/}
\item Calcalist Tech --- Tessl \$125M Funding: \url{https://www.calcalistech.com/ctechnews/article/b1pnxlmg1x}
\item Tessl Product Hunt Launch: \url{https://www.producthunt.com/products/tessl}
\end{itemize}

\newpage

% ================================================================
% STAGE 3 — MISSION PACKAGE
% ================================================================
\stageheader{STAGE 3 --- MISSION PACKAGE}{tertiary}

\section{Milestone Specification}

\subsection{Mission Objective}

Produce a fully-informed PR review decision for gsd-skill-creator PR \#28, including: per-file accept/reject/adapt decisions with technical rationale, a draft PR review comment, GSD-native skill improvements inspired by the PR, and a contribution governance policy for future external skill modifications.

\subsection{Architecture Overview}

\begin{asciibox}
\begin{verbatim}
Wave 0: Foundation                    Wave 1: Analysis (Parallel)
+---------------------------+         +---------------------------+
| 0A: Decision schema       |         | 1A: DACP diff analysis    |
| 0B: Spec compliance       |    +--->| 1B: MFE diff analysis     |
|     checklist              |    |    | 1C: Action security audit |
+-------------+-------------+    |    +-------------+-------------+
              |                   |                  |
              +-------------------+                  |
                                                     v
Wave 2: Synthesis                     Wave 3: Delivery
+---------------------------+         +---------------------------+
| 2A: Decision matrix        |         | 3A: PR review comment     |
| 2B: GSD-native improvement |         | 3B: Contribution policy   |
|     recommendations        |         | 3C: Verification          |
+---------------------------+         +---------------------------+
\end{verbatim}
\end{asciibox}

\subsection{System Layers}

\begin{enumerate}[leftmargin=2em]
\item \textbf{Decision Layer} --- Accept/reject/adapt matrix with technical rationale per file
\item \textbf{Communication Layer} --- PR review comment respecting both technical precision and professional courtesy
\item \textbf{Improvement Layer} --- GSD-native improvements that capture the PR's valid insights without external dependencies
\item \textbf{Governance Layer} --- Contribution policy preventing future context-blind skill modifications
\end{enumerate}

\subsection{Deliverables}

\begin{center}
\rowcolors{2}{rowgray}{white}
\begin{tabularx}{\textwidth}{C{0.5cm}L{4cm}XL{2.5cm}}
\toprule
\rowcolor{primary}
\tableheader{\#} & \tableheader{Deliverable} & \tableheader{Acceptance Criteria} & \tableheader{Component Spec} \\
\midrule
1 & Decision matrix & All 6 files with accept/reject/adapt + rationale & Wave 2A \\
2 & PR review comment & Technically precise, respectful, actionable & Wave 3A \\
3 & GSD improvement list & $\geq$3 concrete improvements with specs & Wave 2B \\
4 & DACP skill patch & Accepted changes applied, bugs fixed & Wave 1A \\
5 & MFE description patches & ``Use when'' clauses added in GSD style & Wave 1B \\
6 & Contribution policy & Skill modification governance rules & Wave 3B \\
\bottomrule
\end{tabularx}
\end{center}

\subsection{Component Breakdown}

\begin{center}
\rowcolors{2}{rowgray}{white}
\begin{tabularx}{\textwidth}{L{3.5cm}L{2.5cm}C{1.5cm}C{2cm}}
\toprule
\rowcolor{primary}
\tableheader{Component} & \tableheader{Dependencies} & \tableheader{Model} & \tableheader{Est. Tokens} \\
\midrule
0A: Decision schema & None & Haiku & $\sim$3K \\
0B: Spec checklist & 0A & Sonnet & $\sim$5K \\
1A: DACP diff analysis & 0A, 0B & Sonnet & $\sim$8K \\
1B: MFE diff analysis & 0A, 0B & Opus & $\sim$12K \\
1C: Action security audit & 0A & Sonnet & $\sim$6K \\
2A: Decision matrix & 1A, 1B, 1C & Opus & $\sim$10K \\
2B: GSD improvements & 2A & Opus & $\sim$8K \\
3A: PR review comment & 2A & Opus & $\sim$5K \\
3B: Contribution policy & 2A, 1C & Sonnet & $\sim$5K \\
3C: Verification & All & Sonnet & $\sim$4K \\
\bottomrule
\end{tabularx}
\end{center}

\subsection{Model Assignment Rationale}

\begin{center}
\rowcolors{2}{rowgray}{white}
\begin{tabularx}{\textwidth}{L{2cm}C{1.5cm}X}
\toprule
\rowcolor{primary}
\tableheader{Model} & \tableheader{Allocation} & \tableheader{Assigned To} \\
\midrule
Opus & $\sim$35\% & MFE analysis (domain judgment), decision matrix, GSD improvements, PR comment \\
Sonnet & $\sim$55\% & DACP analysis, action audit, spec checklist, contribution policy, verification \\
Haiku & $\sim$10\% & Decision schema scaffolding \\
\bottomrule
\end{tabularx}
\end{center}

\section{Wave Execution Plan}

\subsection{Wave Summary}

\begin{center}
\rowcolors{2}{rowgray}{white}
\begin{tabularx}{\textwidth}{C{1.5cm}L{3cm}C{2cm}C{2cm}C{2cm}}
\toprule
\rowcolor{primary}
\tableheader{Wave} & \tableheader{Name} & \tableheader{Execution} & \tableheader{Est. Tokens} & \tableheader{Windows} \\
\midrule
0 & Foundation & Sequential & $\sim$8K & 1 \\
1 & Analysis & Parallel (3 tracks) & $\sim$26K & 2 \\
2 & Synthesis & Sequential & $\sim$18K & 2 \\
3 & Delivery & Sequential & $\sim$14K & 2--3 \\
\bottomrule
\end{tabularx}
\end{center}

\subsection{Wave 0: Foundation}

\begin{center}
\rowcolors{2}{rowgray}{white}
\begin{tabularx}{\textwidth}{C{1cm}XL{3cm}C{1.5cm}C{1.5cm}}
\toprule
\rowcolor{primary}
\tableheader{Task} & \tableheader{Description} & \tableheader{Produces} & \tableheader{Model} & \tableheader{Tokens} \\
\midrule
0A & Define per-file decision schema (accept/reject/adapt fields) & \texttt{types.yaml} & Haiku & $\sim$3K \\
0B & Generate spec compliance checklist for all 5 affected skills & Checklist doc & Sonnet & $\sim$5K \\
\bottomrule
\end{tabularx}
\end{center}

\subsection{Wave 1: Parallel Analysis}

\textbf{Track A --- DACP Analysis:}

\begin{center}
\rowcolors{2}{rowgray}{white}
\begin{tabularx}{\textwidth}{C{1cm}XL{3cm}C{1.5cm}C{1.5cm}}
\toprule
\rowcolor{primary}
\tableheader{Task} & \tableheader{Description} & \tableheader{Produces} & \tableheader{Model} & \tableheader{Tokens} \\
\midrule
1A.1 & Line-by-line diff analysis of DACP skill & Diff assessment & Sonnet & $\sim$5K \\
1A.2 & Generate corrected DACP skill (accepted changes + bug fixes) & Patched SKILL.md & Sonnet & $\sim$3K \\
\bottomrule
\end{tabularx}
\end{center}

\textbf{Track B --- MFE Analysis:}

\begin{center}
\rowcolors{2}{rowgray}{white}
\begin{tabularx}{\textwidth}{C{1cm}XL{3cm}C{1.5cm}C{1.5cm}}
\toprule
\rowcolor{primary}
\tableheader{Task} & \tableheader{Description} & \tableheader{Produces} & \tableheader{Model} & \tableheader{Tokens} \\
\midrule
1B.1 & Analyze all 4 MFE description changes & Description assessment & Opus & $\sim$4K \\
1B.2 & Evaluate summary flattening impact on downstream consumers & Impact report & Opus & $\sim$4K \\
1B.3 & Assess workflow section appropriateness for reference skills & Workflow assessment & Opus & $\sim$4K \\
\bottomrule
\end{tabularx}
\end{center}

\textbf{Track C --- Governance:}

\begin{center}
\rowcolors{2}{rowgray}{white}
\begin{tabularx}{\textwidth}{C{1cm}XL{3cm}C{1.5cm}C{1.5cm}}
\toprule
\rowcolor{primary}
\tableheader{Task} & \tableheader{Description} & \tableheader{Produces} & \tableheader{Model} & \tableheader{Tokens} \\
\midrule
1C.1 & GitHub Action security and data flow audit & Security report & Sonnet & $\sim$4K \\
1C.2 & Review Tessl CLI network behavior documentation & Network analysis & Sonnet & $\sim$2K \\
\bottomrule
\end{tabularx}
\end{center}

\subsection{Wave 2: Synthesis}

\begin{center}
\rowcolors{2}{rowgray}{white}
\begin{tabularx}{\textwidth}{C{1cm}XL{3cm}C{1.5cm}C{1.5cm}}
\toprule
\rowcolor{primary}
\tableheader{Task} & \tableheader{Description} & \tableheader{Produces} & \tableheader{Model} & \tableheader{Tokens} \\
\midrule
2A & Compile per-file decision matrix with rationale & Decision matrix & Opus & $\sim$10K \\
2B & Identify and spec GSD-native improvements & Improvement specs & Opus & $\sim$8K \\
\bottomrule
\end{tabularx}
\end{center}

\subsection{Wave 3: Delivery}

\begin{center}
\rowcolors{2}{rowgray}{white}
\begin{tabularx}{\textwidth}{C{1cm}XL{3cm}C{1.5cm}C{1.5cm}}
\toprule
\rowcolor{primary}
\tableheader{Task} & \tableheader{Description} & \tableheader{Produces} & \tableheader{Model} & \tableheader{Tokens} \\
\midrule
3A & Draft PR review comment & Comment text & Opus & $\sim$5K \\
3B & Write contribution governance policy & Policy doc & Sonnet & $\sim$5K \\
3C & Run verification matrix & Verification report & Sonnet & $\sim$4K \\
\bottomrule
\end{tabularx}
\end{center}

\subsection{Cache Optimization Strategy}

\begin{itemize}[leftmargin=2em]
\item \textbf{Wave 0 types cached for Wave 1:} Decision schema loaded once, used by all three parallel tracks
\item \textbf{PR diff cached across tracks:} The six-file diff is shared context for tracks A, B, and C
\item \textbf{Spec checklist cached:} Generated in 0B, consumed by 1A, 1B, and 2A
\item \textbf{Pre-loaded project knowledge:} gsd-skill-creator-analysis.md and gsd-mathematical-foundations-conversation.md loaded for MFE analysis context
\end{itemize}

\subsection{Token Budget}

\begin{center}
\rowcolors{2}{rowgray}{white}
\begin{tabularx}{\textwidth}{L{3cm}C{2cm}C{2.5cm}C{2cm}}
\toprule
\rowcolor{primary}
\tableheader{Wave} & \tableheader{Est. Tokens} & \tableheader{Context Windows} & \tableheader{Sessions} \\
\midrule
Wave 0 & $\sim$8K & 1 & 1 \\
Wave 1 & $\sim$26K & 2 (parallel) & 1 \\
Wave 2 & $\sim$18K & 2 & 1 \\
Wave 3 & $\sim$14K & 2--3 & 1 \\
\midrule
\textbf{Total} & \textbf{$\sim$66K} & \textbf{$\sim$8} & \textbf{$\sim$3} \\
\bottomrule
\end{tabularx}
\end{center}

\subsection{Dependency Graph}

\begin{asciibox}
\begin{verbatim}
0A (schema) ----+---> 1A (DACP) --------+
                |                        |
0B (checklist) -+---> 1B (MFE) ---------+--> 2A (decisions) --> 3A (comment)
                |                        |                  |
                +---> 1C (governance) ---+                  +--> 3B (policy)
                                                            |
                                                            +--> 3C (verify)

                                         2B (improvements) --^
\end{verbatim}
\end{asciibox}

\section{Test Plan}

\textbf{Total Tests:} 32 | \textbf{Safety-Critical:} 6 | \textbf{Target Coverage:} 85\%+

\subsection{Test Categories}

\begin{center}
\rowcolors{2}{rowgray}{white}
\begin{tabularx}{\textwidth}{L{3cm}C{1.5cm}C{2cm}C{2cm}}
\toprule
\rowcolor{primary}
\tableheader{Category} & \tableheader{Count} & \tableheader{Priority} & \tableheader{Failure Action} \\
\midrule
Safety-critical & 6 & Mandatory & BLOCK \\
Core functionality & 14 & Required & BLOCK \\
Integration & 8 & Required & BLOCK \\
Edge cases & 4 & Best-effort & LOG \\
\bottomrule
\end{tabularx}
\end{center}

\subsection{Safety-Critical Tests}

\begin{center}
\rowcolors{2}{rowgray}{white}
\begin{tabularx}{\textwidth}{C{1cm}XX}
\toprule
\rowcolor{primary}
\tableheader{ID} & \tableheader{Verifies} & \tableheader{Expected Behavior} \\
\midrule
SC-01 & MFE metadata preservation & YAML frontmatter metadata (chapter, plane-position, radius, primitives) identical before and after any accepted changes \\
SC-02 & DACP Safety Invariant visibility & SAFE-01 (no auto-execute) remains prominently visible, not buried in description text \\
SC-03 & No external dependencies introduced & No new CI workflows, action references, or tool installations in merged result \\
SC-04 & Terminology preservation & ``Primitives'' not renamed to ``Concepts'' in any accepted changes \\
SC-05 & PR comment professional tone & Review comment contains appreciation, specific technical rationale, and actionable suggestions \\
SC-06 & No skill content sent externally & Verify all improvements generated locally without Tessl CLI involvement \\
\bottomrule
\end{tabularx}
\end{center}

\subsection{Verification Matrix}

\begin{center}
\rowcolors{2}{rowgray}{white}
\begin{tabularx}{\textwidth}{C{0.5cm}XL{2.5cm}C{1.5cm}}
\toprule
\rowcolor{primary}
\tableheader{\#} & \tableheader{Success Criterion} & \tableheader{Test ID(s)} & \tableheader{Status} \\
\midrule
1 & Accept/reject/adapt decision per file & CF-01, CF-02 & Pending \\
2 & MFE metadata preservation & SC-01 & Pending \\
3 & DACP restructuring evaluated & CF-03, CF-04 & Pending \\
4 & Action security risks quantified & CF-05, SC-03 & Pending \\
5 & PR comment drafted & SC-05, CF-06 & Pending \\
6 & $\geq$3 GSD-native improvements & CF-07, CF-08 & Pending \\
7 & Spec compliance gaps cataloged & CF-09, CF-10 & Pending \\
8 & Contribution policy documented & CF-11, SC-03 & Pending \\
9 & Workflow pattern evaluated & CF-12, IN-01 & Pending \\
10 & Tessl scoring methodology understood & CF-13, CF-14 & Pending \\
\bottomrule
\end{tabularx}
\end{center}

\section{Mission Crew Manifest}

\begin{center}
\rowcolors{2}{rowgray}{white}
\begin{tabularx}{\textwidth}{L{2.5cm}L{2cm}X}
\toprule
\rowcolor{primary}
\tableheader{Role} & \tableheader{Agent} & \tableheader{Responsibility} \\
\midrule
FLIGHT & Orchestrator & Merge decision authority; wave go/no-go gates \\
PLAN & Planner & Wave decomposition; parallel track optimization \\
CRAFT-SPEC & Opus & Agent Skills spec compliance analysis \\
CRAFT-MATH & Opus & MFE domain integrity guardian \\
EXEC-1 & Sonnet & DACP diff analysis and patching \\
EXEC-2 & Sonnet & Governance and security analysis \\
VERIFY & Sonnet & Cross-module consistency checks \\
INTEG & Sonnet & Decision matrix compilation \\
CAPCOM & HITL & Human approval at merge decision point \\
BUDGET & Haiku & Token budget tracking \\
SURGEON & Sonnet & Analysis quality monitoring \\
LOG & Haiku & Commit messages; audit trail \\
\bottomrule
\end{tabularx}
\end{center}

\section{Execution Summary}

\begin{center}
\rowcolors{2}{rowgray}{white}
\begin{tabularx}{\textwidth}{L{5cm}X}
\toprule
\rowcolor{primary}
\tableheader{Metric} & \tableheader{Value} \\
\midrule
Total estimated tokens & $\sim$66K \\
Context windows & $\sim$8 \\
Sessions & $\sim$3 \\
Parallel tracks (Wave 1) & 3 \\
Activation profile & Squadron (12 roles) \\
Critical path & Wave 0 $\rightarrow$ Wave 1B (MFE, highest judgment) $\rightarrow$ Wave 2A $\rightarrow$ Wave 3A \\
Human gates & 1 (merge decision after Wave 2A) \\
\bottomrule
\end{tabularx}
\end{center}

\begin{quotebox}
\emph{``The Amiga's genius wasn't raw power --- it was architectural leverage. A 7MHz 68000 with 512K of RAM had no business doing what the Amiga did.''}

This PR is an Amiga moment in miniature. The power of GSD's mathematical foundations skills isn't in their Tessl scores --- it's in the architectural relationships between them. A tool that measures structural compliance against a generic spec is measuring clock speed when it should be measuring DMA channel throughput. Take what's genuinely useful. Protect what makes the architecture sing. And build the governance so the next contribution arrives into a system that knows how to evaluate itself.
\end{quotebox}

\end{document}
